\newcommand{\engFunc}[1][P]{\measure{F}{}{Q,\tilde{#1}}}
\newcommand{\infR}{\inf_{\inSet{R}}}
\newcommand{\supR}{\sup_{\inSet{R}}}

Suppose we want to approximate a central distribution of a hierarchy of distributions~\set{h}.
To do so, we need a measure that quantifies the quality of an approximation.
A natural choice is the central divergence introduced in Definition~\ref{def:central-divergence}.
Since central distributions are intended to serve as representatives of hierarchies of distributions,
central divergence can be viewed as a measure of how well a probability distribution represents a hierarchy.

At first glance, computing this quantity appears to require explicit knowledge of the central distributions of~\set{h}.
However, it turns out that the central divergence admits an upper bound that can be estimated without explicitly
knowing the central distribution.

In variational settings, it is often more convenient to work with unnormalized measures rather than normalized
distributions.
Therefore, we state the next theorem in terms of unnormalized measures.

\begin{theorem}\label{th:central-approx}
    Let \setQ be a finite set of probability distributions over a discrete random variable~$X$.
    Let $\tilde{P}$ be a positive measure over $X$,
    and let $P$ be the corresponding normalized distribution.
    Let $R$ be an arbitrary convex combination of the elements of $\set{R} \subseteq \setQ$.
    Then
    \[
        \myDiv[P][Q][] \leq \supR \engFunc - \inf_{\inSet{Q}} \engFunc + \kl[R][P],
    \]
    where
    \[
        \engFunc = \entropy{Q}{X} + \Ex{Q}{\ln \tilde{P}(X)}.
    \]
\end{theorem}

\begin{proof}
    By Corollary~\ref{cor:cdiv-lambda-bound}, we have
    \[
        \myDiv[P][Q][] \leq \sFunc[][\setQ] - g(\cDist).
    \]
    where \cDist is the \cd.

    Since $R$ is a convex combination of the elements of \set{R}, Lemma~\ref{lm:kl-bound} implies
    \[
        g(R) \geq \infR \kl - \kl[R][P].
    \]
    By Theorem~\ref{th:C-maximizes-g}, \cDist is a maximizer of $g$, and therefore
    \[
        g(\cDist) \geq g(R).
    \]
    Combining these three inequalities yields
    \begin{equation} \label{eq:cdiv-upper-bound}
        \myDiv[P][Q][]  \leq \sFunc[][\setQ] - \infR \kl + \kl[R][P].
    \end{equation}

    Let $\tilde{P}(X) = Z P(X)$, where $Z > 0$ is a normalization constant.
    Then for any~$Q$,
    \[
        \kl = \ln Z - \engFunc.
    \]
    Since $\ln Z$ does not depend on $Q$, we have
    \begin{align*}
        \infR \kl &= \ln Z - \supR \engFunc \\
        \sup_{\inSet{Q}} \kl &= \ln Z - \inf_{\inSet{Q}} \engFunc,
    \end{align*}
    and the result follows from Equation~\ref{eq:cdiv-upper-bound}.
\end{proof}

By Theorem~\ref{th:central-approx}, it suffices to construct a positive measure $\tilde P$ for which $\engFunc$ is
approximately constant over \setQ.
We may choose this constant to be zero and seek
\[
    \engFunc \approx 0.
\]
In addition, the normalized distribution $P$ should be close to a convex combination of the elements of \setQ.

Since central distributions are defined through marginal distributions, their construction may be viewed as
combining prior specification and Bayesian inference within a single optimization framework.
Given $n$ observations, one may perform inference using the central distribution of the corresponding finite
sequence~$S_n$, without separately specifying a prior distribution.

For sequences of random variables, however, central distributions are defined on the
underlying infinite sequence.
Consequently, the central distribution of a finite initial segment $S_n$ should be viewed as an approximation to the
corresponding marginal of the infinite-sequence central distribution.
The question then becomes how accurately this approximation can be constructed.

One possible approach is to assume that $\ln \tilde P(X)$ can be represented by a linear combination of a
collection of basis functions $b_i$, analogous to the log-density representation used in exponential families:
\[
    \ln \tilde P(x) = -\sum_i \alpha_i b_i(x).
\]
Substituting this into $\measure{F}{}{}$ gives
\[
    \engFunc = \entropy{Q}{X} - \sum_i \alpha_i \Ex{Q}{b_i(X)}.
\]
Therefore, enforcing $\engFunc \approx 0$ amounts to matching the entropy function.

Suppose that \setQ is indexed by a parameter \param.
Let \set{S} be a finite set of samples of parameter values.
We may then choose the coefficients $\alpha_i$ so that for every \inSet[\paramVal]{S}
\[
    \sum_i \alpha_i \Ex{}{b_i(X) \mid \paramVal} \approx \entropy{Q}{X \mid \paramVal}.
\]
The resulting probability distribution is given by
\[
    P(x) \propto \exp \left\{-\sum_i \alpha_i b_i(x) \right\}.
\]

The quality of the approximation depends on how well $\tilde P$ can be represented as a nonnegative linear
combination of the distributions $Q(X \mid \paramVal)$.
This can be approximately assessed using the sample set \set{S} by finding nonnegative weights $w_i$ such that
\[
    \tilde P(x) \approx \sum_i w_i Q(x \mid \paramVal_i)
\]
where $\inSet[\paramVal_i]{S}$ and $w_i\ge 0$.
The following example illustrates the approximation procedure in a simple binomial setting.

\begin{example}
    \newcommand{\Fdelta}{$\mathbb{F}$-Delta}
\newcommand{\PredErr}{Pred-Err}

As an illustration, consider a binomial model with parameter $\theta$.
Let $S_n = (X_1, X_2, \ldots, X_n)$ denote a sample of $n$ independent observations,
where $X_i \sim \operatorname{Bernoulli}(\theta)$.

We use the basis functions
\[
    b_k(s_n) \coloneqq \left(\frac1n \sum_{i=1}^n x_i \right)^k.
\]
These basis functions form polynomial features of the sufficient statistic.

For a binomial model with success probability $\theta$, the entropy of $S_n$ is
\[
    \entropy{Q}{S_n \mid \paramVal} = -n \left(\paramVal \ln \paramVal + (1 - \paramVal) \ln (1 - \paramVal)  \right).
\]

We consider the case $n=7$ and approximate $\cDist(S_7)$ using the procedure described above.
The parameter space is represented by a sample set \set{S} consisting of 30 values of $\theta$.

When the approximation lies close to a marginal likelihood within the binomial family,
the bound in Theorem~\ref{th:central-approx} can be estimated by
\[
    \max_{\inSet[\paramVal]{S}} \engFunc - \min_{\inSet[\paramVal]{S}} \engFunc,
\]
which we denote by \Fdelta.

To investigate the effect of the number of basis functions on the quality of the approximation, we vary the number
of basis functions from 4 to 14.
The corresponding values of \Fdelta\ are shown in Figure~\ref{fig:error_comparision}.

\begin{figure}[t]
    \centering
    \begin{tikzpicture}
        \begin{axis}[
            globalstyle,
            legend pos=north east,
            xlabel={Number of Basis},
            xmin=4, xmax=14,
            ymin=0.0
        %xtick={3,...,14},
        ]

            \addplot+[mark=*, blue] coordinates {
                (4, 0.622701) (5, 0.192069) (6, 0.192069) (7, 0.076361) (8, 0.076361) (9, 0.028878) (10, 0.028878) (11, 0.011228) (12, 0.015049) (13, 0.012935) (14, 0.014727)
            };
            \addlegendentry{\Fdelta}

        \end{axis}
    \end{tikzpicture}
    \begin{tikzpicture}
        \begin{axis}[
            globalstyle,
            legend pos=north east,
            xlabel={Number of Basis},
            xmin=4, xmax=14,
            ymin=0.0
        %xtick={3,...,14},
        ]

            \addplot+[mark=*, blue] coordinates {
                (4, 0.033566) (5, 0.012962) (6, 0.012962) (7, 0.006028) (8, 0.006028) (9, 0.000312) (10, 0.000312) (11, 0.000675) (12, 0.000083) (13, 0.000426) (14, 0.000171)
            };
            \addlegendentry{\PredErr}

            \addplot[very thick, dashed, red] coordinates {(4, 0.01785714) (14, 0.01785714)};
            \addlegendentry{Uni-Diff}
        \end{axis}
    \end{tikzpicture}
    \caption{
        Approximation quality as a function of the number of basis functions.
        Left: \Fdelta, the bound derived from Theorem~\ref{th:central-approx}.
        Right: \PredErr\, the absolute difference between the estimated predictive probability and the
        central predictive probability.
        The dotted line (Uni-Diff) indicates the corresponding error obtained using a uniform prior.
    }
    \label{fig:error_comparision}
\end{figure}


Next, we use the approximation to estimate the central predictive probability $C(X_7\mid S_6)$.
Here, \cDist denotes the marginal distribution induced by the central prior $\theta^{-\hf}(1 - \theta)^{-\hf}$.
We consider an observed sequence $S_6$ containing two successes and four failures.
Under \cDist, the predictive probability of success is $\frac{5}{14}$.
Figure~\ref{fig:error_comparision} reports the absolute difference between the estimated predictive probability and
the central predictive probability, denoted by \PredErr.
For comparison, the absolute difference of the predictive probability obtained from the
uniform prior is also shown as the dotted line (Uni-Diff).

The results show that as the number of basis functions increases, both \Fdelta\ and \PredErr\ decrease substantially.
This suggests that improvements in the approximation bound are accompanied by improvements in the resulting
predictive distribution.
Moreover, with a sufficient number of basis functions, the predictive estimate becomes considerably closer to the
central posterior than the estimate obtained from the uniform prior.


\end{example}